\ documentclass [12 pt ,a4 paper ]{ article }
\ usepackage [ utf 8]{ inputenc }
\ usepackage [ portuguese ]{ babel }
\ usepackage { booktabs , longtable , graphicx , hyperref }
\ usepackage { geometry }
\ geometry { margin =2.5 cm }
% METADADOS -- preencher obrigatoriamente
\ newcommand {\ projetotitulo }{ Projeto Conecte - Implementação de Tecnologias Digitais e Metodologias Ativas na Educação }
\ newcommand {\ grupoid }{ Grupo X}
\ newcommand {\ sprintnum }{ N}
\ newcommand {\ sprintperiodo }{ dd / mm a dd / mm / aaaa }
\ newcommand {\ scrummaster }{ Nome do SM }
\ newcommand {\ gitrepo }{ https :// github . com / grupo - x/ projeto }
\ newcommand {\ gitcommit }{ abc 1234} % hash do commit desta entrega
\ begin { document }
\ begin { center }
{\ Large \ bfseries Relatorio de Sprint \ sprintnum }\\
{\ large \ projetotitulo }\\
\ grupoid {} | \ sprintperiodo {} | SM : \ scrummaster
\ end { center }
\ section { Sprint Backlog }
% Tabela : ID | Historia | Estimativa | Responsavel | Status
\ begin { longtable }{ cllcc }
\ toprule
ID & Historia & Estimativa & Responsavel & Status \\
\ midrule
% ... linhas aqui
\ bottomrule
\ end { longtable }
\ section { Burndown Chart }
\ begin { figure }[ h !]
\ centering
\ includegraphics [ width =0.7\ textwidth ]{ burndown _ sprint \
sprintnum . png }
\ caption { Burndown -- Sprint \ sprintnum }
\ end { figure }
% Analise do burndown ( minimo 3 linhas ) :
\ section { Artefatos Produzidos }
% Descrever com evidencias
\ section { Atualizacao com o Parceiro }
% Data , formato e feedback
\ section { Participacao do Calouro }
% Atividade realizada ou justificativa de ausencia
\ section { Impedimentos e Desvios }
\ section { Retrospectiva }
\ subsection *{ O que funcionou bem }
\ subsection *{ O que precisa melhorar }
\ subsection *{ Acao concreta para a proxima sprint }
\ section { Planejamento da Proxima Sprint }
\ end { document }